% Auto-generated by python/build_tables.py
\begin{table*}[!htbp]
\centering
\begin{threeparttable}
\caption{Baseline architecture comparison.}
\label{tab:baseline-architecture-results}
\small
\begin{tabular}{p{0.30\textwidth}rrrrr}
\toprule
\textbf{Architecture} & \textbf{Mean $T$} & \textbf{Median $T$} & \textbf{$P_{90}$} & \textbf{$P_{95}$} & \textbf{Convergence} \\
\midrule
Baseline modular network & 7,435.8 & 7,048.0 & 9,925.6 & 11,016.9 & 100\% \\
Random bridging & 7,280.2 & 7,008.4 & 9,653.6 & 10,592.6 & 100\% \\
Boundary spanning & 4,576.0 & 4,513.5 & 5,444.7 & 5,743.5 & 100\% \\
\bottomrule
\end{tabular}
\begin{tablenotes}
\footnotesize
\item Notes: Values are graph-level summaries of first-passage time to relational coordination readiness. Each architecture uses $NG=50$ graph realizations and $NT=50$ trajectories per graph. Random bridging does not clearly outperform the modular baseline, whereas boundary spanning substantially reduces readiness time and upper-tail delay.
\end{tablenotes}
\end{threeparttable}
\end{table*}
